\paragraph{Letra A}

Neste caso a planta escolhida deve ser instável, escolhemos arbitrariamente uma função de transferência da forma $G_(s) = 2/(s-10)$, com o polo localizado no semi-plano direito. Os parâmetros a serem estimados são a posição do polo no semi plano direito e o ganho da planta: $a=10; b=2$. Introduzimos novamente um esquema de normalização e fazemos as estimativas. Notamos que em ambos os casos a saída do modelo da planta não converge exatamente para a saída da planta, porém o erro de saída normalizado $\varepsilon_o$ converge, e a diferença entre os logaritmos naturais da saída do modelo e da planta nunca se afasta de uma determinada faixa de valores, que fica firmemente atrelada à escolha das estimativas iniciais. 


\includegraphics{./pic/Q3letraAmontagem.png}
 
Podemos observar o comportamento da estimativa de a no gráfico a seguir:

\includegraphics{./pic/Q3letraAteta1.png}

E da estimativa de b no gráfico a seguir:

\includegraphics{./pic/Q3LetraAteta2.png} 

E o erro normalizado convergindo para 0:

\includegraphics{./pic/Q3letraAe0.png}

E a diferença entre os logaritmos naturais da saída da planta e do modelo

\includegraphics{./pic/Q3letraALxLxc.png}
 

\paragraph{Letra B}

Aqui usamos um sinal senoidal para fazer a estimativa e obtivemos resultados similares.


\includegraphics{./pic/Q3letraBmontagem.png}
 
Podemos observar o comportamento da estimativa de a no gráfico a seguir:

\includegraphics{./pic/Q3letraBteta1.png}

E da estimativa de b no gráfico a seguir:

\includegraphics{./pic/Q3LetraBteta2.png} 

E o erro normalizado convergindo para 0:

\includegraphics{./pic/Q3letraBe0.png}

E a diferença entre os logaritmos naturais da saída da planta e do modelo

\includegraphics{./pic/Q3letraBLxLxc.png}
 
 